\documentclass[11pt]{amsart}
\usepackage[T1]{fontenc}
\usepackage{lmodern,amsmath,amssymb,mathtools,microtype,booktabs}
\usepackage[margin=1.15in]{geometry}
\usepackage[hidelinks]{hyperref}
\newtheorem{theorem}{Theorem}[section]
\newtheorem{proposition}[theorem]{Proposition}
\newtheorem{lemma}[theorem]{Lemma}
\newtheorem{corollary}[theorem]{Corollary}
\theoremstyle{definition}\newtheorem{definition}[theorem]{Definition}
\theoremstyle{remark}\newtheorem{remark}[theorem]{Remark}
\newcommand{\R}{\mathbb R}
\newcommand{\E}{\mathbb E}
\newcommand{\Prob}{\mathbb P}
\newcommand{\dd}{\,\mathrm d}
\newcommand{\one}{\mathbf 1}
\DeclareMathOperator{\diag}{diag}
\DeclareMathOperator{\arcosh}{arcosh}
\numberwithin{equation}{section}
\title[Intrinsic marked boundary laws]{Intrinsic marked boundary laws and rigidity of periodic dispersing billiards}
\author{Qian Qi}
\date{September 28, 2026}
\subjclass[2020]{37D50, 37C30, 53A04, 70H09, 62G05}
\keywords{Dispersing billiards, intrinsic arclength observations, boundary laws,
analytic rigidity, periodic contact networks, inverse problems}
\hypersetup{pdftitle={Intrinsic marked boundary laws and rigidity of periodic dispersing billiards},pdfauthor={Qian Qi}}
\begin{document}
\begin{abstract}
We prove rigidity theorems from intrinsic arclength-marked, near-onset
collision laws in periodic dispersing billiards. At an isolated labelled
normal channel, the observation consists of the two endpoint-arclength
subprobability laws, considered up to one simultaneous reversal. Their
onset, masses, signed first moments and two leading quadratic moments
recover the gap, the separate curvatures, the free area and every contact
jet. No numerical contact frame or leading geometry is supplied. The
inverse has explicit nonsingular even and odd blocks at every degree,
including equal curvatures. For analytic boundaries, the observation
determines the participating pair up to Euclidean isometry. A connected
finite channel network, with intrinsic arclength registration and
rank-two lattice-cycle marking, determines the entire periodic table
and its marked lattice. We also describe the positive-dimensional
finite-jet fibers of the unmarked count law, separating this question
from equality of exact analytic count germs. The Cartesian signed-moment
inverse, physical realization, finite-window experiments and smooth
relative boundary theory are retained. No finite-horizon hypothesis is
imposed.
\end{abstract}
\maketitle
\section{Intrinsic observations and the rigidity theorem}
\label{sec:introduction}
Near the first possible occurrence of a collision itinerary, its
probability integrates both the action of the broken ray and the physical
flux through a collision section. We study the geometric information in
this law. The central observation is intrinsic: endpoint positions are
recorded by distance along the reflecting boundary, not by a supplied
Cartesian contact chart. A single reversal of all the arclength signs
is immaterial. The resulting inverse recovers the leading geometry as
well as the higher asymmetric jets.

There are two distinct steps from a local law to table rigidity.
First, a finite compression of the marked law determines a participating
pair of analytic obstacles. Second, intrinsic registration of the contact
points identifies the pairwise reconstructions on a connected channel
network. Lattice-cycle marking then determines the periods. The second
step requires genuine inter-channel information; it does not assert that
an unvisited obstacle is determined by one short itinerary.

\subsection{Geometric and observational conventions}
Consider finitely many labelled compact strictly convex obstacles modulo
a full-rank lattice in $\R^2$. Their boundaries are smooth with positive
curvature, their closures are disjoint, and the free cell area is $A>0$.
Initial conditions have normalized phase volume
$\dd q\dd\vartheta/(2\pi A)$ and speed one. A selected channel is an
isolated unobstructed common normal segment between two designated
obstacle copies. Its endpoints are labelled $0,1$. Fixed sufficiently
small facing patches isolate this channel; clearance and positive
separation guarantee a common small time collar.

The two endpoint origins are part of the channel marking. Their ambient
coordinates, the normal direction, the gap, the curvatures, and the area
are not supplied numerically. Choose either orientation of the plane
solely to describe the observation. Let $n$ point from contact $0$ to
contact $1$ and let $t$ be its positive quarter-turn. At both endpoints
use signed arclength increasing in tangent direction $t$. Equivalently,
one may use counterclockwise arclength around each obstacle and reverse
its sign at contact $1$. This prescription involves no numerical
Cartesian chart. Reversing the plane orientation reverses both signs.
The same choice is used for both itineraries and all windows.

For a time window $T$, let $E_b(T)$ be the event of exactly the three
successive impacts $b,1-b,b$ in the selected patches. On success record
the first and last signed boundary arclengths $(s_-,s_+)$ from contact
$b$; failures have a separate symbol. Define the subprobability measure
\begin{equation}\label{eq:intrinsicmeasure}
 \mu_{b,T}(B)=\Prob\{E_b(T),\ (s_-,s_+)\in B\},\qquad B\subset\R^2.
\end{equation}
The intrinsic observation is the orbit of the entire family
$(\mu_{0,T},\mu_{1,T})_T$ under
$\iota(s_-,s_+)=(-s_-,-s_+)$. It is \emph{not} the average of a
measure with its reflection, nor a separate sign choice at each time.
In particular, the observation retains odd information while supplying
no preferred orientation. It uses a boundary-arclength sensor and is
strictly richer than a collision count.

Write $T_0$ for the local onset of positive mass and $d=T-T_0$.
From any one representative of the observation extract
\begin{align}
 P_b(d)&=\mu_{b,T_0+d}(\R^2),&
 M_b(d)&=\int(s_-+s_+)\dd\mu_{b,T_0+d},\label{eq:arcdatum}\\
 D_b(d)&=\int(s_--s_+)^2\dd\mu_{b,T_0+d},&
 \beta_b&=\lim_{d\downarrow0}\frac{D_b(d)}{dP_b(d)}.\label{eq:betadef}
\end{align}
Thus the proof below uses only $T_0$, two mass germs, two first-moment
germs, and two numbers $\beta_b$, not the whole marked distribution.
The moments are unconditional; no observation or preparation is discarded
on failure. The quotient of these quantities by their single common sign
is a canonical finite compression of \eqref{eq:intrinsicmeasure}.

\begin{theorem}[Intrinsic local and registered global rigidity]
\label{thm:intrinsic}
The following assertions hold for the observation above.

\smallskip\noindent
\textup{(i)} Its onset and quadratic moments recover the leading geometry:
\begin{equation}\label{eq:calibrationmain}
 g=\frac{T_0}{2},\qquad
 \kappa_b=\frac{2}{3\beta_b}-\frac1g,\qquad
 A=\frac1{4\alpha\sqrt{c_0c_1(c_0c_1-1)}},
\end{equation}
where $c_b=1+g\kappa_b$ and
$\alpha=\lim_{d\downarrow0}P_b(d)/d^2$ is independent of $b$.

\smallskip\noindent
\textup{(ii)} At every fixed order $K\ge3$, these leading data and the
coefficients of $P_b/(\alpha d^2)$ and $M_b/(\alpha d^2)$ through the
orders specified in Theorem~\ref{thm:main} determine the contact $K$-jets
up to simultaneous transverse reflection. On either coherent orientation
cover the inverse is analytic and locally bi-Lipschitz at finite order.
Its highest-jet blocks are exactly \eqref{eq:evenblock} and
\eqref{eq:oddblock}; arbitrary lower odd jets and equal curvatures are
allowed. For connected closed real-analytic obstacle boundaries, the
observation determines the entire participating pair, in its relative
position, up to one Euclidean isometry.

\smallskip\noindent
\textup{(iii)} Suppose every obstacle is visited by a connected finite
registered normal-channel network as defined in
Section~\ref{sec:global}. If the marked lattice displacements of its
cycles span $\R^2$, its intrinsic observations and registration determine
the whole labelled periodic table and its marked lattice up to one
Euclidean isometry. No obstacle image or numerical contact frame is
supplied.
\end{theorem}

The main theorem concerns exact germs. The finite-order stability in
(ii) is not stability of analytic continuation, is not uniform in $K$,
and does not turn an onset limit into a finite-sample measurement.
Part (iii) includes explicit intrinsic registration and lattice labels;
these are additional data, not consequences of a single channel law.
At reflection-fixed jets the unframed space is a finite quotient rather
than a smooth coordinate manifold. The asserted analytic coordinates
are on a coherent orientation cover.

\subsection{Cartesian coordinates used in the proof}
After the gap has been identified, put the two contacts at $(0,0)$ and
$(g,0)$. The facing boundaries have inward height functions
\begin{equation}\label{eq:graphs}
 \psi_b(y)=\frac{\kappa_by^2}{2}
       +\sum_{r\ge3}\frac{q_{b,r}}{r!}y^r,\qquad \kappa_b>0.
\end{equation}
The actual positions are $(-\psi_0(y),y)$ and
$(g+\psi_1(y),y)$. For smooth boundaries the series denotes jets only.
The corresponding signed arclength is
\begin{equation}\label{eq:arcmap}
 \mathfrak a_b(y)=\int_0^y\sqrt{1+\psi_b'(x)^2}\dd x
       =y+\frac{\kappa_b^2y^3}{6}+O(y^4).
\end{equation}
Hence $M_b$ is the expectation of $\mathfrak a_b(u)+\mathfrak a_b(v)$ on success.
The auxiliary Cartesian moment and normalized coefficients are
\begin{gather}
 R_b(d)=\E[(u+v)\one_{E_b(2g+d)}],\qquad
 E_{2,b}(d):=E_b(2g+d),\label{eq:rawdata}\\
 z=c_0c_1-1,\quad \gamma=\arcosh\sqrt{c_0c_1},\quad
 \alpha=\frac1{2A\sinh(2\gamma)},\label{eq:leading}\\
 G_b(d)=\frac{P_b(d)}{\alpha d^2}=1+\sum_{k\ge1}\xi_{b,k}d^k,
 \quad J_b(d)=\frac{R_b(d)}{\alpha d^2}=\sum_{k\ge1}\eta_{b,k}d^k.
 \label{eq:coefficients}
\end{gather}
The Cartesian moment is used to compute universal blocks; it is not
silently substituted for the intrinsic observable. Lemma~\ref{lem:filter}
proves the precise transfer to arclength.

\begin{theorem}[Cartesian signed-moment component]\label{thm:main}
At supplied labelled gap, separate curvatures and common transverse
orientation the following assertions hold.
\textup{(i)} The common leading probability coefficient is $\alpha$ and
determines $A$. \textup{(ii)} For
$n_e=\lfloor K/2\rfloor-1$, $n_o=\lfloor(K-1)/2\rfloor$,
the coefficient vectors $(\xi_{b,k})_{1\le k\le n_e}$ and
$(\eta_{b,k})_{1\le k\le n_o}$, $b=0,1$, are analytic triangular
coordinates for $(q_{b,r})_{3\le r\le K}$, locally bi-Lipschitz at
fixed order on compact positive leading-geometry boxes.
\textup{(iii)} For analytic contacts the four raw germs $P_b,R_b$
determine the participating analytic boundary images in these frames.
\textup{(iv)} Analytic periodic physical families realize all the displayed
jet directions, at fixed area or with area free. Respectively $2K-4$ or
$2K-3$ positive-window scalar means are local bi-Lipschitz coordinates
on these supplied families. At fixed area their selected leading
count and endpoint-metric hierarchy is fixed at every flight number.
\end{theorem}

We retain the full Cartesian proofs, physical realization and charged
binary experiment in Sections~\ref{sec:local}, \ref{sec:inverse} and
\ref{sec:physical}. The standalone filtration below supplies the
arbitrary-lower-jet justification, and Lemma~\ref{lem:analyticidentity}
supplies the full continuation statement used in part (iii).

\subsection{Relation to inverse billiard problems}
B\'alint, De Simoi, Kaloshin and Leguil~\cite{BDKL} recover period-two
curvatures and periodic Lyapunov exponents from marked lengths of open
dispersing billiards. De Simoi, Kaloshin and Leguil~\cite{DSKL} obtain
analytic table determination under symmetry and genericity assumptions.
Osterman~\cite{Osterman} relates marked lengths to homoclinic analytic
conjugacy and determines a third scatterer from two known scatterers and
the marked spectrum. Finamore and Leguil~\cite{FL} prove isometric
rigidity from an enriched marked length spectrum for finite-horizon
Sinai billiards. These are global dynamical observations; none is
identified here with a near-onset arclength law.

The distinction of mechanism also matters. Koval~\cite{Koval} studies
formal billiard reconstruction and Gevrey divergence with odd derivatives
as parameters. Our highest-jet calculation is a physical residual-flux
integral, not an assertion that every formal inverse series converges.
De Simoi, Kaloshin and Wei~\cite{DSKW} prove deformational dynamical
spectral rigidity near a circle in a symmetric convex class. Their
periodic-length observation and focusing geometry differ from those
used here. We make no assertion of exhaustive priority from this
comparison.

The present result removes supplied leading geometry by a quadratic
moment calibration, replaces Cartesian marks by intrinsic metric marks,
and joins local analytic determinations by finite registration. It does
not claim count-only asymmetric rigidity. Section~\ref{sec:countfibers}
describes the finite and formal count-data fibers, while distinguishing
these from the exact analytic count-germ question. The smooth relative
boundary law and its function-valued Volterra and Abel theory retain their
own information category in Section~\ref{sec:relative} and Supplement S.

\section{The local two-flight law}
\label{sec:local}
Put $o=1-b$. With longitudinal displacement taken into the corresponding
obstacles, one flight has length
\[
 \ell_{bo}(u,w)=
 \sqrt{\{g+\psi_b(u)+\psi_o(w)\}^2+(u-w)^2}.
\]
The broken action is
$\ell_{bo}(u,w)+\ell_{ob}(w,v)$. Its $w$ derivative vanishes at
$(u,w,v)=(0,0,0)$, with second $w$ derivative $2c_o/g>0$.
The implicit-function theorem gives a unique stationary point
$w=w_b(u,v)$ on a fixed endpoint box. The segment clearance and
reflection law ensure that it is the actual local three-impact path.
Let
\[
 W_b(u,v)=\ell_{bo}(u,w_b)+\ell_{ob}(w_b,v),\qquad
 E_b=W_b-2g.
\]

\begin{lemma}[Quadratic action and physical density]\label{lem:flux}
Writing $y=(u,v)^t$, one has
\begin{equation}\label{eq:quad}
 w_b=\frac{u+v}{2c_o}+O(|y|^2),\qquad
 E_b(y)=E_{b,0}(y)+O(|y|^3),\qquad
 E_{b,0}(y)=\tfrac12 y^tH_by,
\end{equation}
where
\begin{equation}\label{eq:H}
 H_b=\frac1g
 \begin{pmatrix}c_b-(2c_o)^{-1}&-(2c_o)^{-1}\\
                 -(2c_o)^{-1}&c_b-(2c_o)^{-1}\end{pmatrix},
 \quad a_b=\sqrt{\det H_b}=\frac1g\sqrt{\frac{c_bz}{c_o}}.
\end{equation}
The normalized twist
\begin{equation}\label{eq:twist}
 b_b(y)=\frac{-W_{b,uv}(y)}{d_b^0},\qquad
 d_b^0=\frac1{2gc_o}=\frac{a_b}{\sinh(2\gamma)}
\end{equation}
is positive on a smaller box and equals one at the origin. On a
sufficiently small fixed offset collar,
\begin{align}
 G_b(d)&=\frac{a_b}{\pi d^2}
             \int(d-E_b(y))_+b_b(y)\dd y,\label{eq:Gintegral}\\
 J_b(d)&=\frac{a_b}{\pi d^2}
             \int(u+v)(d-E_b(y))_+b_b(y)\dd y.\label{eq:Jintegral}
\end{align}
Both functions extend smoothly to zero, with $G_b(0)=1$ and
$J_b(0)=0$. The constructions and every fixed derivative are uniform
on compact local smooth families with the necessary higher bounds.
\end{lemma}
\begin{proof}
The quadratic broken action is
\[
 \frac{c_bu^2+2c_ow^2+c_bv^2-2w(u+v)}{2g}.
\]
Eliminating $w$ proves \eqref{eq:quad}--\eqref{eq:H}.
Both eigenvalues of $H_b$ are positive, since $c_bc_o>1$.
The mixed derivative at zero gives \eqref{eq:twist}.
In the coordinates of a collision section, Liouville volume has density
$\dd u\dd p\dd r$, where $p$ is momentum conjugate to $u$ and $r$
is time before the first impact. The generating-function identity
$p=-W_{b,u}$ changes its density to
$(-W_{b,uv})\dd u\dd v\dd r$ on this dispersing branch.
The time constraint is $0<r<d-E_b(u,v)$.

Choose the collar smaller than a common positive lower bound on free
flight lengths. Then the initial residual interval does not cross a
preceding collision, and the time remaining after the third collision
cannot reach a fourth one. The strict local minimum of $E_b$ places
its entire active sublevel inside the fixed facing patches when the
collar is small. Every parametrized path is therefore a physical first-hit
path in the specified event; conversely every such near-onset path is
in this chart. Division by $2\pi A$ and integration in $r$ give the raw
versions of \eqref{eq:Gintegral} and \eqref{eq:Jintegral}.
Multiplication by $2A\sinh(2\gamma)/d^2$ proves the displayed formulas.
The same calculation holds with the bounded mark $u+v$, without
changing the measure or conditioning on success.

For completeness, let $y=\Phi_b(x)$ be a local smooth Morse chart,
with $E_b(\Phi_b(x))=|x|^2/2$ and $D\Phi_b(0)=H_b^{-1/2}$.
After $x=\sqrt d\,X$, division of either integral by $d^2$ leaves
an integral on the fixed disk $|X|^2<2$ with weight $1-|X|^2/2$.
For any smooth amplitude its Taylor expansion in $h=\sqrt d$ has only
even integrated powers: the integral at $-h$ equals that at $h$ by
$X\mapsto-X$. Taylor's formula, to arbitrarily high fixed order,
therefore gives a smooth function of $d\ge0$, including parameter
derivatives. For analytic graphs the same argument applies to a
convergent local series. The constant amplitude has integral
\[
 I_0(d):=\int(d-E_{b,0})_+\dd y=\frac{\pi d^2}{a_b}.
\]
This gives $G_b(0)=1$. The marked amplitude vanishes at the origin,
so $J_b(0)=0$ and $J_b(d)=O(d)$.
\end{proof}

The first part of Theorem~\ref{thm:main} follows immediately. Notice
that area is extracted from a raw exact germ before the normalization
in \eqref{eq:coefficients}; no unknown area is inserted into a finite
observation.

\section{Recovering the leading geometry from intrinsic marks}
\label{sec:calibration}
All arclength variables in this section refer to one representative of
\eqref{eq:intrinsicmeasure}. The quantities used for calibration are
unchanged when that representative is reversed.

\begin{proposition}[Onset and quadratic calibration]\label{prop:calibration}
The local onset is $T_0=2g$. The conditional rescaled endpoint law
converges to the probability measure on $\R^2$ with density
\begin{equation}\label{eq:ellipselimit}
 \frac{a_b}{\pi}(1-\tfrac12Y^tH_bY)_+\dd Y.
\end{equation}
In particular, its covariance is $H_b^{-1}/3$ and
\begin{equation}\label{eq:betavalue}
 \beta_b=\frac{2g}{3c_b},\qquad
 \beta_b^+:=\lim_{d\downarrow0}
 \frac{\int(s_-+s_+)^2\dd\mu_{b,2g+d}}{dP_b(d)}
 =\frac{2g}{3(c_b-c_{1-b}^{-1})}.
\end{equation}
Consequently \eqref{eq:calibrationmain} holds. The plus moment is an
optional compatibility check; it is not needed by the inverse.
\end{proposition}
\begin{proof}
The reduced action has a strict minimum $2g$ at the normal orbit by
Lemma~\ref{lem:flux}. On a sufficiently small fixed patch no competing
minimum occurs. There is no event for $T\le2g$, whereas the positive
twist and the open residual interval give positive mass for every
sufficiently small $T>2g$. This proves the onset statement.

Let $y=\sqrt d\,Y$ in the physical integral. Its active set stays in a
bounded ellipse, because $E_b(y)\ge c|y|^2$ near the minimum. Taylor's
formula gives $E_b(\sqrt dY)/d\to E_{b,0}(Y)$,
$b_b(\sqrt dY)\to1$ and
$\mathfrak a_b(\sqrt dY_i)/\sqrt d\to Y_i$ uniformly there. The constant $a_b$ in
\eqref{eq:ellipselimit} is $\sqrt{\det H_b}$, as in \eqref{eq:H}.
After division by $P_b(d)$, dominated convergence gives
\eqref{eq:ellipselimit}, including its bounded quadratic moments.

Under $X=H_b^{1/2}Y$, the normalizing integral is
$\int_{|X|<\sqrt2}(1-|X|^2/2)\dd X=\pi$.
Its second moment in either coordinate is $\pi/3$, and its mixed
moment is zero. Thus the covariance is $H_b^{-1}/3$.
The vector $(1,-1)^t$ is an eigenvector of $H_b$ with eigenvalue
$c_b/g$, and $(1,1)^t$ has eigenvalue $(c_b-c_{1-b}^{-1})/g$.
This proves \eqref{eq:betavalue} without knowing either curvature in
advance. Solving its first formula gives
$\kappa_b=2/(3\beta_b)-1/g$. Finally
$\sinh(2\gamma)=2\sqrt{c_0c_1(c_0c_1-1)}$ and the leading physical
probability coefficient give the area formula.
\end{proof}

Only intrinsic boundary distances enter this computation. The formula is
not an inversion of the support of an embedded endpoint cloud: no
longitudinal boundary coordinate is observed. The two curvatures are
separated by the two itineraries even at equal curvature.

\begin{corollary}[Leading-coordinate Jacobian]\label{cor:leadingchart}
On $g,\kappa_0,\kappa_1,A>0$, the map
$(g,\kappa_0,\kappa_1,A)\mapsto(T_0,\beta_0,\beta_1,\alpha)$
has an analytic inverse. In this order its determinant is
\begin{equation}\label{eq:leadingdet}
 -\frac{8\alpha g^4}{9A c_0^2c_1^2}\ne0.
\end{equation}
Together with the higher-jet blocks it gives an analytic $2K$-dimensional
finite-jet coordinate map on a chosen orientation cover.
\end{corollary}
\begin{proof}
After the first row, eliminate the $g$ derivatives in the two
$\beta$ rows. Their remaining diagonal derivatives are
$-2g^2/(3c_b^2)$. The last diagonal derivative is
$\partial_A\alpha=-\alpha/A$. Multiplication with
$\partial_gT_0=2$ gives \eqref{eq:leadingdet}. The number of higher
coordinates is $2(K-2)$.
\end{proof}

\section{A filtered inverse for intrinsic arclength}
\label{sec:filtered}
The following lemma makes the highest-jet calculation independent of
any symmetry assumption on the lower jets. It also records precisely
why a nonlinear arclength mark has the same leading blocks as its
Cartesian linearization. We use the arclength map
$\mathfrak a_b(y)$ of \eqref{eq:arcmap}.

\begin{lemma}[Residual-flux filtration]\label{lem:filter}
Fix the leading geometry and an integer $r\ge3$. Let $\theta$ be one
contact derivative of degree $r$, all lower derivatives being arbitrary
and fixed. Suppose the reduced action, normalized twist and a mark have
variations
\begin{align*}
 \partial_\theta E(y)&=e_r(y)+O(|y|^{r+1}),&
 \partial_\theta b(y)&=b_{r-2}(y)+O(|y|^{r-1}),\\
 f(y)&=f_1(y)+O(|y|^2),&
 \partial_\theta f(y)&=O(|y|^{r+1}),
\end{align*}
where the subscripts denote homogeneous degrees and $f_1$ is a fixed
linear form. The leading geometry is independent of $\theta$. Normalize
$\int(d-E)_+b\dd y$ and $\int f(d-E)_+b\dd y$ by $I_0(d)$.
For even $r=2m$, the derivative of the count coefficient of $d^{m-1}$
is the coefficient obtained using only $E_0$, $e_r$, and $b_{r-2}$.
For odd $r=2m+1$, the derivative of the marked coefficient of $d^m$
is obtained using only those terms and $f_1$.
If these homogeneous variations depend only on the leading geometry,
so do the highest-jet derivatives, at every choice of lower asymmetric
jets. Both coefficient functionals are affine in their highest jets.
\end{lemma}
\begin{proof}
A Morse chart makes the normalized integrals smooth functions of $d$
and of finite-dimensional smooth parameters, as in
Lemma~\ref{lem:flux}. Differentiation before scaling gives the exact
first-variation formula
\begin{equation}\label{eq:firstvariationfilter}
 \partial_\theta\int f(d-E)_+b\dd y
 =-\int_{E<d}f b\,\partial_\theta E\dd y
  +\int(d-E)_+\{f\partial_\theta b+b\partial_\theta f\}\dd y.
\end{equation}
There is no boundary term, since the residual factor vanishes on
$E=d$. The same formula with $f=1$ applies to counts.
Put $d=h^2$, $y=hY$. After normalization the action and twist
variations have order $h^{r-2}$; a linear mark adds one power of $h$.
The variation of the mark itself starts at $h^{r+1}$. More explicitly,
the coefficient of $h^{r-1}$ for the marked variation is
\begin{equation}\label{eq:filteredcoefficient}
 \frac{a_b}{\pi}\left\{-\int_{E_0<1}f_1e_r\dd Y
              +\int(1-E_0)_+f_1b_{r-2}\dd Y\right\}.
\end{equation}
For counts remove $f_1$ and replace $r-1$ by $r-2$.

Here the error estimate is uniform on compact finite-jet neighborhoods.
The rescaled energy is $E_0+O(h)$ in $C^1$, and the level $E_0=1$
is regular and compact. Its domain differs from the quadratic domain
by a region of volume $O(h)$. The twist is $1+O(h)$ and the rescaled
mark is $f_1+O(h)$. Substitution into
\eqref{eq:firstvariationfilter} therefore changes the displayed leading
term by at least one additional power of $h$. This explicitly accounts
for the moving domain; no nonlinear Morse-map term has been omitted.
The fixed-domain expression is even in $h$, so odd integrated powers
vanish and the coefficient of the indicated even power is well defined.

A derivative of degree higher than $2m$ first enters the count variation
above its target weight $2m-2$. A derivative of degree higher than
$2m+1$ first enters the marked variation above $2m$. Thus only the
stated finite jets occur. The stationary equation is solved degree by
degree by division by $2c_{1-b}/g$; its remaining operations and the
finite quadratic moments are analytic in the leading geometry and
polynomial in the finite higher jets. Since the displayed derivatives
are independent of the higher jets, integration in the highest-jet
variables proves affinity, not merely a linearization at even contacts.
\end{proof}

In the billiard problem the homogeneous variations are
\eqref{eq:variationE}--\eqref{eq:variationb}, proved by stationary
action in Section~\ref{sec:inverse}. They satisfy the lemma at arbitrary
lower jets. For the intrinsic mark there is one further check:
\begin{equation}\label{eq:arcvariation}
 \partial_{q_{b,r}}\mathfrak a_b(y)
 =\int_0^y\frac{\psi_b'(x)}{\sqrt{1+\psi_b'(x)^2}}
                    \frac{x^{r-1}}{(r-1)!}\dd x
 =\frac{\kappa_b y^{r+1}}{(r+1)(r-1)!}+O(y^{r+2}).
\end{equation}
The opposite contact does not enter this endpoint mark. The dependence
of arclength on the unknown graph therefore appears \emph{two weights
later} than the odd action contribution. This is the required
noncircularity in passing from Cartesian to intrinsic observations.

\begin{proposition}[Arclength coefficient inverse]\label{prop:arcinverse}
Write
\begin{equation}\label{eq:arcseries}
 \widetilde J_b(d)=\frac{M_b(d)}{\alpha d^2}
               =\sum_{k\ge1}\widetilde\eta_{b,k}d^k.
\end{equation}
Then $D_{Q_{2m+1}}\widetilde\eta_m$ is exactly
\eqref{eq:oddblock}. The even count block is \eqref{eq:evenblock}.
In graph-degree order
$\widetilde\eta_1,\xi_1,\widetilde\eta_2,\xi_2,\ldots$,
these are triangular coordinates on every finite contact-jet space.
\end{proposition}
\begin{proof}
The mark is $\mathfrak a_b(u)+\mathfrak a_b(v)$, with linear part
$u+v$. Equations \eqref{eq:arcmap} and \eqref{eq:arcvariation} satisfy
Lemma~\ref{lem:filter}. Its highest derivative is therefore precisely
the quadratic moment computed in Proposition~\ref{prop:blocks}.
The remainder at each degree is an analytic function only of previously
recovered jets and of the now calibrated leading geometry. Subtracting
it and solving the displayed invertible block gives the recursion
\eqref{eq:recursion}, with the arclength remainder in place of the
Cartesian one. The strict determinant inequalities
\eqref{eq:evenpositive}--\eqref{eq:oddpositive} require neither unequal
curvatures nor a nonzero higher derivative.
\end{proof}

\begin{lemma}[Identity of connected analytic boundary images]
\label{lem:analyticidentity}
Let $C$ and $C'$ be compact connected embedded real-analytic
one-dimensional submanifolds of $\R^2$ without boundary. If their
images share a nonempty open regular arc, then $C=C'$.
\end{lemma}
\begin{proof}
Let $S\subset C$ be the points at which the two submanifold germs
coincide. This set is nonempty and relatively open. If $x_j\in S$
converges to $x\in C$, compactness of $C'$ implies $x\in C'$.
Continuity of tangent lines at the common arcs implies that the tangent
lines of $C$ and $C'$ at $x$ agree. In a sufficiently small box both
are graphs of real-analytic functions over this common tangent line.
If $x_j=x$ for some $j$, the germ equality is already known.
Otherwise their difference has zeros accumulating at the interior point
$x$ and the one-variable identity theorem makes the two graphs equal.
Thus $S$ is closed in $C$. Connectedness gives $S=C$, so $C\subset C'$.
This inclusion is relatively open by the germ equality and closed by
compactness. Connectedness of $C'$ gives equality. The proof is about
embedded images and hence introduces no choice of continued
parametrization and no monodromy assumption.
\end{proof}

\begin{proof}[Proof of Theorem~\ref{thm:intrinsic}(i)--(ii)]
Proposition~\ref{prop:calibration} gives (i). On one orientation cover,
Proposition~\ref{prop:arcinverse} recovers all finite jets, and
Corollary~\ref{cor:leadingchart} includes the unknown leading geometry.
At fixed order the finite denominators are bounded away from zero on
compact positive leading-geometry sets; bounded derivatives of the map
and its inverse give the local bi-Lipschitz assertion.

For analytic contacts equality of the recovered jets gives equality of
the local graphs. Place the two contacts at $(0,0)$ and $(g,0)$ and
apply Lemma~\ref{lem:analyticidentity} separately to the two closed
boundaries. Their relative position is fixed by this same channel frame.
Changing the observation representative negates $M_b$, leaves $P_b$
and $\beta_b$ unchanged, and sends $q_{b,r}$ to $(-1)^rq_{b,r}$.
The recovery is equivariant under this transformation. Its two
representatives therefore differ by a common reflection, not by
independent reflections of the two obstacles. Translation and rotation
of the chosen numerical channel frame were arbitrary to begin with.
This proves the intrinsic isometry statement.
\end{proof}

\section{The all-degree inverse}
\label{sec:inverse}
For the rest of the algebra keep $g,\kappa_0,\kappa_1$ fixed and put
\begin{equation}\label{eq:constants}
\begin{gathered}
 L_b=\frac{g}{2c_bz},\qquad \mathsf R=1+2z,\qquad V_m=1+2mz,\\
 k_m=\frac4{2^m(m+1)(m!)^2},\qquad
 C_m=\frac{2^{2-m}}{(m+1)(m+2)(m!)^2}.
\end{gathered}
\end{equation}
Here $\mathsf R$ is not the raw first moment $R_b$.
Use column vectors $Q_r=(q_{0,r},q_{1,r})^t$,
$\xi_k=(\xi_{0,k},\xi_{1,k})^t$, and
$\eta_k=(\eta_{0,k},\eta_{1,k})^t$.

\begin{proposition}[Successive even and odd blocks]\label{prop:blocks}
For every $m\ge2$, $\xi_{m-1}$ depends only on $Q_3,\ldots,Q_{2m}$,
is affine in $Q_{2m}$, and
\begin{equation}\label{eq:evenblock}
 D_{Q_{2m}}\xi_{m-1}
 =-\begin{pmatrix}\mathsf R^m&V_m\\V_m&\mathsf R^m\end{pmatrix}
                      \diag(k_mL_0^m,k_mL_1^m).
\end{equation}
For every $m\ge1$, $\eta_m$ depends only on $Q_3,\ldots,Q_{2m+1}$,
is affine in $Q_{2m+1}$, and
\begin{equation}\label{eq:oddblock}
 D_{Q_{2m+1}}\eta_m
 =-\begin{pmatrix}c_1\mathsf R^m&V_m\\
                   V_m&c_0\mathsf R^m\end{pmatrix}
              \diag(2c_0C_mL_0^{m+1},2c_1C_mL_1^{m+1}).
\end{equation}
No evenness assumption is needed for either assertion. The symmetric
matrix in each formula is positive definite. In particular,
\begin{align}
 \mathsf R^m-V_m&=\sum_{r=2}^m\binom mr(2z)^r>0
                          &&(m\ge2),\label{eq:evenpositive}\\
 c_0c_1\mathsf R^{2m}-V_m^2
   &\ge z(1+2mz)^2>0 &&(m\ge1).\label{eq:oddpositive}
\end{align}
\end{proposition}

\subsection{Finite-jet dependence}
Vary one graph derivative of order $r\ge3$, keeping all lower graph
jets fixed. Stationarity of the intermediate point gives the leading
variation of the reduced action without differentiating that point.
With $w_0=(u+v)/(2c_o)$ one obtains
\begin{align}
 \partial_{q_{b,r}}E_b&=\frac{u^r+v^r}{r!}+O(|y|^{r+1}),
       &\partial_{q_{o,r}}E_b&=\frac{2w_0^r}{r!}+O(|y|^{r+1}),
                                                \label{eq:variationE}\\
 \partial_{q_{b,r}}b_b&=O(|y|^{r-1}),
       &\partial_{q_{o,r}}b_b&=-\frac{g}{c_o}
                    \frac{w_0^{r-2}}{(r-2)!}+O(|y|^{r-1}).
                                                \label{eq:variationb}
\end{align}
Indeed the graph perturbation changes the longitudinal flight length
to first order with coefficient one at the normal orbit; each endpoint
is counted once and the intermediate contact twice. The correction to
$w_0$ has order two, so it changes no displayed homogeneous term.
Twice differentiating the action and using $d_b^0=(2gc_o)^{-1}$
proves the twist formulas. This is where the contribution of the
opposite contact to the density enters.

Here and below the remainder statements mean Taylor remainders with
the finite number of derivatives needed at the chosen order. The
coefficient assertions can be justified on the fixed Morse disk in
Lemma~\ref{lem:flux}. Equivalently, at the scaling $y=hY$, $d=h^2$,
a graph jet of degree $r$ first enters the rescaled action and twist
with weight $h^{r-2}$. The mark adds a factor $h$. Every nonquadratic
unvaried correction adds at least one further power of $h$. Odd powers
integrate to zero after the fixed-domain construction.
Consequently $[d^{m-1}]G_b$ uses only graph jets through degree $2m$,
whereas $[d^m]J_b$ uses only those through degree $2m+1$.
At their respective highest degrees the displayed homogeneous
variations are integrated against the quadratic action and constant
twist. Products with any lower nonlinear correction have higher weight.
The resulting highest-jet derivative is independent of all jets of
order at least three. This proves affinity, not just linearization at
a symmetric reference contact.

To see explicitly that these are analytic finite-jet maps, differentiate
the stationary equation at the origin degree by degree. Its highest
unknown derivative is always multiplied by $2c_o/g>0$. The remaining
operations are finite sums, products, and division by positive
leading-geometry quantities. Integration of the finitely many terms on
the quadratic disk preserves analytic dependence. Neither analyticity
nor evenness of the original smooth graph is needed for this argument.

\subsection{Quadratic moments}
For a linear form $\ell(y)$ set
$\nu(\ell)=\ell^tH_b^{-1}\ell$, and for two forms let
$\operatorname{cov}(t,\ell)=t^tH_b^{-1}\ell$.
Inverting \eqref{eq:H} gives
\begin{equation}\label{eq:covariances}
 H_b^{-1}=L_b\begin{pmatrix}\mathsf R&1\\1&\mathsf R\end{pmatrix},
 \qquad \nu(u)=L_b\mathsf R,\quad \nu(w_0)=L_o.
\end{equation}
With $t=u+v$,
\begin{equation}\label{eq:tcov}
 \operatorname{cov}(t,u)=\operatorname{cov}(t,v)
        =2c_bc_oL_b,\qquad t=2c_ow_0.
\end{equation}
Rotation after the transformation $y=H_b^{-1/2}x$ gives
\begin{align}
 \frac1{I_0(d)}\int_{E_{b,0}<d}\ell^{2r}\dd y
 &=\frac{2(2r)!\nu(\ell)^r}{2^r(r!)^2(r+1)}d^{r-1},
                                          \label{eq:diskmoment}\\
 \frac1{I_0(d)}\int(d-E_{b,0})_+\ell^{2r}\dd y
 &=\frac{2\binom{2r}{r}\nu(\ell)^r}{2^r(r+1)(r+2)}d^r.
                                          \label{eq:weightedmoment}
\end{align}
For example, the angular mean is
$\binom{2r}{r}/4^r$ and the unweighted radial integral is
$(2d)^{r+1}/(2r+2)$. For either radial weight and every $r\ge0$,
\begin{equation}\label{eq:crossmoment}
 \int t\ell^{2r+1}\,\mathrm{weight}
   =\frac{\operatorname{cov}(t,\ell)}{\nu(\ell)}
             \int\ell^{2r+2}\,\mathrm{weight}.
\end{equation}
To prove this identity, decompose $t$ in the transformed coordinates
into its component along $\ell$ and the orthogonal component. The latter
integrates to zero by reflection. All these formulas are exact.

\subsection{Computation of the blocks}
The first action variation in either integral is the negative integral
of the action perturbation over $E_{b,0}<d$, with the appropriate mark.
There is no extra moving-boundary term because the residual weight
vanishes on the boundary.
For $r=2m$, the two endpoint powers in \eqref{eq:variationE} give
$-k_mL_b^m\mathsf R^m$. The opposite-contact action gives
$-k_mL_o^m$. Its twist variation, by \eqref{eq:weightedmoment}, gives
$-k_mL_o^m(2mz)$, since $g/(c_oL_o)=2z$.
The sum is the first row of \eqref{eq:evenblock}; reversing the
contacts gives the other row. This calculation remains valid when
lower odd jets are nonzero, by the weight argument above.

For $r=2m+1$, use the mark $t=u+v$. Equations
\eqref{eq:diskmoment} and \eqref{eq:crossmoment}, with
$r$ there replaced by $m+1$, yield the own-contact action contribution
\begin{equation}\label{eq:ownodd}
 -C_m\operatorname{cov}(t,u)\nu(u)^m
       =-2c_bc_oC_mL_b^{m+1}\mathsf R^m.
\end{equation}
The opposite-contact action contributes $-2c_oC_mL_o^{m+1}$.
For its twist, use $t=2c_ow_0$ and \eqref{eq:weightedmoment}:
\begin{align*}
 &-\frac{g}{c_o(2m-1)!I_0(d)}
       \int t\,w_0^{2m-1}(d-E_{b,0})_+\dd y\\
 &\hspace{1cm}=-2m g C_mL_o^m d^m
       =-4mc_ozC_mL_o^{m+1}d^m.
\end{align*}
Thus the opposite-contact coefficient is
$-2c_oC_mL_o^{m+1}(1+2mz)$. Together with \eqref{eq:ownodd}
this proves \eqref{eq:oddblock}.

For the even symmetric matrix, \eqref{eq:evenpositive} and the
positive sum of the two entries prove positivity. For the odd one,
Bernoulli's inequality gives $\mathsf R^m\ge1+2mz$ for all
$m\ge1$. Since $c_0c_1=1+z$, its determinant is at least
$z(1+2mz)^2$. Its diagonal entries are positive. This proves
Proposition~\ref{prop:blocks} at every order.

At $g=1$, $c_0=c_1=2$, the first odd block is
\begin{equation}\label{eq:cubicexample}
 D_{Q_3}\eta_1=-\frac7{108}
                  \begin{pmatrix}2&1\\1&2\end{pmatrix}.
\end{equation}
Its nonzero antisymmetric eigenvalue is already visible at cubic
order. The first even block is
$D_{Q_4}\xi_1=-\frac1{1728}\left(\begin{smallmatrix}49&13\\13&49\end{smallmatrix}\right)$.
These examples are normalizations of the proof, not a substitute for it.

\subsection{Recursive recovery and analytic determination}
Order the data by graph degree as
\[
 \eta_1,\ \xi_1,\ \eta_2,\ \xi_2,\ldots .
\]
At degree $r$, let $Y_r$ be the corresponding two-component datum,
let $B_r$ be its block in Proposition~\ref{prop:blocks}, and let
$F_r(Q_3,\ldots,Q_{r-1})$ be its value when $Q_r=0$.
Then the inverse is the exact recursion
\begin{equation}\label{eq:recursion}
 Q_r=B_r^{-1}\{Y_r-F_r(Q_3,\ldots,Q_{r-1})\},\qquad r\ge3.
\end{equation}
Every inverse block exists. A fixed finite collection of its denominators
has a positive minimum on a compact positive leading-geometry box.
The recursion is analytic; on sufficiently small compact coordinate
neighborhoods the map and inverse have bounded derivatives. Integration
along segments gives the local bi-Lipschitz conclusion. This proves
part (ii) of Theorem~\ref{thm:main}.

For analytic contacts, the four raw germs first determine $A$ and then
all the coefficients in \eqref{eq:coefficients}. Recursion
\eqref{eq:recursion} identifies every graph derivative. Analyticity
identifies the graph germs. To obtain the boundary images, continue the
common analytic regular curve along arclength. At each point the local
graph identity continues uniquely; compactness and connectedness of
the closed embedded boundaries continue it around the entire curve.
This proves part (iii). The continuation concerns the participating
boundaries, not an obstacle never visited by the selected channel.

\begin{remark}[Why the signed information is essential]
The transformation $\psi_b(y)\mapsto\psi_b(-y)$ at both contacts
sends $(u,w,v)$ to $(-u,-w,-v)$. It leaves $P_b$ unchanged and sends
$R_b$ to $-R_b$. Whenever an odd graph jet is nonzero, these are distinct
oriented contact germs. Thus an orientation-sensitive inverse cannot
be deduced from the unmarked probability germs by simply removing an
evenness hypothesis. The mark in \eqref{eq:rawdata} supplies exactly
the odd information used by \eqref{eq:oddblock}. Even when all $R_b$
vanish, the recursion applies: in conjunction with the probability
germs it forces the successive odd jets to vanish.
\end{remark}

\section{From intrinsic channel laws to the periodic table}
\label{sec:global}
A local contact law does not register different experiments in a common
plane. This section specifies finite intrinsic data that do so. The
registration consists of boundary distances and integer copy labels,
not positions, angles, curvatures, or known obstacle images.

\begin{definition}[Registered normal-channel network]\label{def:network}
Let $V$ index all the obstacle components in one cell. A finite directed
multigraph has an edge $e=(v,w,k_e)$ for each selected isolated clear
normal channel from a chosen copy of obstacle $v$ to the copy of obstacle
$w$ displaced by $Lk_e$, where $k_e\in\mathbb Z^2$ and $L$ is the
unknown lattice basis matrix. Loops and multiple edges are allowed.
The integer labels are part of the observation. The underlying graph
on $V$ is connected, and every vertex has an incident channel.

At each vertex distinguish one incident contact occurrence. For every
other incident occurrence supply its signed arclength distance from that
reference along the same boundary, modulo its perimeter. All distances
and all channel-law representatives use one coherent choice of boundary
orientation, considered modulo its single simultaneous reversal. At the
source endpoint a channel's transverse arclength has that boundary
orientation, and at the target it has the opposite orientation, as in
Section~\ref{sec:introduction}. Contact occurrences, not just edges,
are registered, including the two occurrences of a loop. No independent
sign reversal at a vertex or edge is allowed.

For a closed graph walk the signed sum of its integer edge labels is
its cycle label. The network has \emph{full lattice rank} if these labels
span $\R^2$.
\end{definition}

The orientation in this definition is intrinsic bookkeeping. Equality
of registered data means equality after one common sign choice, not
equality of independently symmetrized edge laws. Perimeters need not
be supplied separately: they are determined when a participating
analytic boundary image has been reconstructed. The registration must
be compatible with actual arclength on that image; the theorem assumes
the data come from physical tables.

\begin{theorem}[Registered periodic rigidity]\label{thm:network}
Let two periodic dispersing tables have connected closed real-analytic
obstacle boundaries. If their registered normal-channel networks have
full lattice rank and their intrinsic channel laws and registration agree,
then one Euclidean isometry takes the whole first labelled periodic table
and its marked lattice to the second. The lattice need not be supplied
numerically. The conclusion uses only the compressed local data in
\eqref{eq:arcdatum}--\eqref{eq:betadef}, their onsets, and the finite
registration of Definition~\ref{def:network}.
\end{theorem}
\begin{proof}
Choose matching representatives of the single orientation orbit. Fix
one vertex and one incident contact, and place that contact and its
outward normal in a numerical reference frame. Theorem~\ref{thm:intrinsic}(i)--(ii),
already proved without using this section, reconstructs its entire
analytic boundary image. Translation and rotation of the initial frame
are the only choices at this step. Arclength registration on this now
known closed curve locates every other incident contact uniquely modulo
the perimeter. In particular, it supplies each contact point and its
outward unit normal without measuring an ambient frame.

Choose a spanning tree rooted at this vertex. Assign integers
$\tau_v\in\mathbb Z^2$ along the tree so that each tree edge has
transformed label
\[
 k'_e=k_e+\tau_v-\tau_w=0.
\]
This only changes the chosen representative copy of each obstacle.
Traverse the tree. When a source contact is placed at $p_{v,e}$ with
outward normal $n_{v,e}$, the partner contact is at
\begin{equation}\label{eq:placepartner}
 p_{v,e}+g_e n_{v,e}.
\end{equation}
Its outward normal is $-n_{v,e}$. The recovered local graph and the
coherent arclength orientation place the partner germ uniquely there;
Lemma~\ref{lem:analyticidentity} places its entire boundary image.
Its registered contacts can then be located. If a tree edge is traversed
backwards, use its target occurrence and the reversed normal segment;
this is the same construction with the endpoints interchanged. Induction
therefore places all representative obstacles in one common plane.
For two tables with matching data, these placed images coincide.

For a remaining edge, both contact occurrences are now known on their
representative images. Its physical displacement must satisfy
\begin{equation}\label{eq:cyclelinear}
 d_e:=p_{v,e}+g_e n_{v,e}-p_{w,e}=L k'_e.
\end{equation}
The non-tree transformed labels generate the real span of all cycle
labels: adjoining a non-tree edge to the tree path gives precisely its
fundamental cycle, with label $k'_e$. Full rank supplies two edges
$e_1,e_2$ whose transformed labels are linearly independent. Hence
\begin{equation}\label{eq:latticerecover}
 L=[d_{e_1}\ d_{e_2}]
                 [k'_{e_1}\ k'_{e_2}]^{-1}.
\end{equation}
The integer matrix need not be unimodular. Its entries are known copy
labels, so invertibility over $\R$ suffices. Every other cycle gives
a consistency equation, not a further choice of periods. Thus the two
tables have the same marked lattice and the same obstacle images after
the initial common placement. Reversing the coherent orientation changes
this placement by one reflection. Restoring the arbitrary initial
translation and rotation proves the isometry assertion.
\end{proof}

This proves Theorem~\ref{thm:intrinsic}(iii). The graph hypothesis is a
finite visibility and registration condition, not a finite-horizon
assumption. It permits the two contacts of a channel to lie on distinct
translates of the same obstacle.

\begin{proposition}[Physical examples and the rank condition]
\label{prop:networkexamples}
There is a nonempty open set of analytic periodic tables admitting the
networks in Theorem~\ref{thm:network}, without reflection symmetry.
If the full-rank condition is omitted, a registered one-channel law can
remain unchanged while the periodic table changes by a nonisometric
lattice shear.
\end{proposition}
\begin{proof}
Start with a unit disk on the lattice $3\mathbb Z\times4\mathbb Z$.
Use the horizontal and vertical normal channels to its neighboring
copies, with labels $(1,0)$ and $(0,1)$. Their gaps are one and two.
The four contact occurrences form the registration on the single
obstacle vertex. The channels are isolated strict minima of the local
distance action, and their segments have positive clearance from all
other obstacles. These properties persist under sufficiently small
analytic $C^2$ perturbations of the support function and of the lattice:
the implicit-function theorem continues the normal contacts, and a
uniform diameter bound reduces the clearance check to finitely many
nearby translates. Arbitrary small nonsymmetric perturbations are
allowed. The two cycle labels retain rank two.

For the last assertion keep the unit disk fixed and use the lattice
basis $(3,0),(s,4)$, with $|s|$ sufficiently small, observing only the
horizontal channel. Its contact arcs, gap, local action and twist remain
unchanged. The covolume is always $12$, so the normalized phase area is
always $12-\pi$. Positive clearance gives a common small observation
collar. The two intrinsic local laws and their contact registration
are therefore exactly the same for all these tables. They do not depend
on the unseen second period. Nevertheless the lattices for $s=0$ and
small $s\ne0$ are not isometric: their shortest vector independent of
$(3,0)$ has length respectively $4$ and $\sqrt{16+s^2}$.
Thus local analytic continuation determines the obstacle in this example,
but not an unobserved lattice direction. Equation~\eqref{eq:latticerecover}
identifies exactly the additional information that removes this freedom.
\end{proof}

\section{The count-only question at finite and formal order}
\label{sec:countfibers}
The signed observation and the unmarked count are different experiments.
A simultaneous transverse reflection leaves both count germs unchanged.
That fact alone does not decide whether exact analytic count germs
determine an asymmetric pair modulo reflection. We give a more precise
finite-order statement, and specify what it does and does not imply at
infinite order.

\begin{proposition}[Finite count fibers]\label{prop:countfibers}
Fix the positive leading geometry and $M\ge2$. The map
\[
 (Q_3,\ldots,Q_{2M})\longmapsto(\xi_1,\ldots,\xi_{M-1})
\]
has analytic local fibers of dimension $2M-2$. The odd jets
$Q_3,Q_5,\ldots,Q_{2M-1}$ are coordinates along these fibers; the
even jets are determined successively by them and the fixed count
coefficients. Near the physical tables of Proposition~\ref{prop:realization},
these fibers occur in actual analytic periodic tables at fixed gap,
curvatures and area. Quotienting by simultaneous reflection does not
remove their positive dimension.
\end{proposition}
\begin{proof}
At stage $m$ the count equation is
\begin{equation}\label{eq:countfiberrecursion}
 \xi_{m-1}=B_{2m}Q_{2m}
             +F_{2m}(Q_3,\ldots,Q_{2m-1}),\qquad 2\le m\le M.
\end{equation}
The matrix $B_{2m}$ is \eqref{eq:evenblock}, and is invertible.
Choose all the odd jets freely in a small neighborhood. At $m=2$
solve \eqref{eq:countfiberrecursion} for $Q_4$; then solve for $Q_6$,
and continue to $Q_{2M}$. This yields an analytic graph over the
$2(M-1)$ odd variables. Conversely every solution is on this graph.
The complete jet-coordinate map in Proposition~\ref{prop:realization}
realizes each nearby solution in an analytic periodic table with the
specified leading geometry and area. A finite reflection group cannot
make a positive-dimensional local analytic fiber discrete: away from
its fixed set its quotient is locally finite-to-one, and arbitrarily
many distinct reflection orbits remain also near the fixed set.
\end{proof}

The coefficient map in this proposition is truncated. Equality of these
$M-1$ coefficients is not equality at finitely many positive windows,
whose higher-order remainders may carry additional information.
The positive-window coordinate theorem for the marked experiment below
is consequently consistent with these count fibers.

\begin{proposition}[Formal fibers and smooth realization]
\label{prop:formalcount}
At fixed leading geometry, any prescribed sequence of formal odd contact
jets determines a unique sequence of formal even jets with prescribed
formal count coefficients, by \eqref{eq:countfiberrecursion}. Near a
fixed strictly dispersing periodic reference geometry, such compatible
formal contact jets can be realized by smooth physical tables with the
same contact locations, gap, curvatures and free area. In particular,
there are pairs not related by common reflection whose normalized count
laws have the same Taylor series at onset.
\end{proposition}
\begin{proof}
The formal assertion follows by induction, since each equation uses only
finitely many earlier jets and the same nonsingular even block. There
is no infinite equation to solve at any induction step.

For the physical assertion one may start with the unit-disk construction
of Proposition~\ref{prop:realization}. The two registered contacts are
distinct points of one boundary modulo the lattice. A prescribed graph
jet there determines, degree by degree, a support-function jet: at each
degree the diagonal support-to-graph derivative is the nonzero value
in \eqref{eq:supportjac}. The leading support value, first derivative
and curvature radius are kept fixed. We recall explicitly the needed
smooth jet extension. In disjoint support-angle neighborhoods of the
two contacts, multiply the desired degree-$r$ Taylor monomial by a
smooth cutoff equal to one near its origin, with radius $\varepsilon_r$.
Choose the radii decreasing so rapidly that the $r$th summand has
$C^k$ norm at most $2^{-r}\epsilon$ for every $k\le r-1$.
This is possible because the degree-$r$ monomial has such norms of
order $\varepsilon_r^{r-k}$, with fixed coefficients at that stage.
The resulting series converges in every fixed $C^k$ norm after omitting
finitely many terms. Since each cutoff is identically one near its
origin, it has exactly the prescribed jets. Beginning at degree three
makes the full perturbation arbitrarily small in $C^2$, regardless of
the later formal coefficients.

Choose, in addition, a nonnegative nonzero smooth support perturbation
$f$ away from both contacts. It changes no contact jet. The free-area
derivative in this direction is
$-\int(h+h'')f\dd\theta<0$. The implicit-function theorem therefore
restores the reference area by a small multiple of $f$. All support
perturbations can be kept sufficiently small in $C^2$ that
$h+h''>0$, embedding, obstacle separation and channel clearance persist.
These are actual smooth periodic billiards. Finite-jet dependence of
the coefficient functionals identifies their count Taylor coefficients
with the prescribed formal ones.

To obtain the last assertion, prescribe the same count series as the
reference table but choose a nonzero cubic jet, or two odd sequences
not related by common negation, before solving the even recursion.
The resulting smooth jets, and hence the boundary germs, are not a
common reflection pair. Their count Taylor series are identical.
\end{proof}

\begin{remark}[Exact analytic count germs]\label{rem:analyticcount}
No convergence assertion for the formal inverse in
Proposition~\ref{prop:formalcount} is used or established. Its smooth
realizations may have count laws differing by a nonzero flat germ.
Thus the proposition neither proves nor disproves determination by
\emph{exact analytic} count germs modulo common reflection. It shows why
finite-order counting and the reflection obstruction alone cannot settle
that question. Establishing convergence of suitable formal fibers, or a
rigidity mechanism that excludes them in the analytic class, would be
needed to decide it. We do not assign either outcome on the basis of a
formal recursion. The intrinsic marked theorem avoids this unresolved
step by determining each odd jet directly with a nonsingular block.
\end{remark}

\section{Physical realization and positive-window observation}
\label{sec:physical}
The coefficient inverse is useful only if its directions occur in actual
tables. We give an analytic realization of every parity, retaining the
contact locations and curvatures. The transverse axis is the same global
axis at the two horizontal contacts; it is not reversed when the
longitudinal contact frame is reflected.

\subsection{Independent all-degree physical jets}
\begin{proposition}[Physical jet coordinates]\label{prop:realization}
Fix $K\ge3$ and two radii $R_0,R_1$ sufficiently close to one.
On the lattice $3\mathbb Z\times4\mathbb Z$ there is an analytic
family of positively curved obstacles for which the horizontal selected
channel has $g=1$, $\kappa_b=R_b^{-1}$ and fixed contact locations.
The variables
\[
 (q_{b,r})_{b=0,1;\,3\le r\le K},\quad A
\]
are local analytic coordinates on this family. Restricting to
$A=12-\pi$ gives a $2(K-2)$-dimensional analytic family with the
selected leading count and endpoint-metric coefficients fixed at every
flight number. Its jet directions of both parities are independent.
\end{proposition}
\begin{proof}
Use one obstacle in a cell and write
\[
 w_0(\theta)=\frac{1+\cos\theta}{2},\quad
 w_1(\theta)=\frac{1-\cos\theta}{2},\quad
 a=\frac{R_0+R_1-2}{4},\quad b=\frac{R_0-R_1}{4}.
\]
Consider the support function
\begin{equation}\label{eq:support}
 \begin{split}
 h(\theta)={}&1+a\sin^2\theta+b\cos\theta\sin^2\theta\\
 &+\sum_{r=3}^K\{s_{0r}w_0(\theta)+s_{1r}w_1(\theta)\}
                                      \sin^r\theta
        +\zeta\sin^{2K+2}\theta.
 \end{split}
\end{equation}
At normals $0,\pi$ its values are one and its first derivatives zero.
Its radii of curvature there are $R_0,R_1$. None of the $s$ or
$\zeta$ variables changes these statements. For small parameters
$h+h''>0$, so the curve is regular, analytic, embedded, and positively
curved. Its two support points remain $(1,0)$ and $(-1,0)$.
The chord to the translate by $(3,0)$ has length one and is normal.

At the unit disk all other translates have a strictly larger gap than
the two horizontal nearest neighbors. This and positive chord clearance
persist for small perturbations. A uniform bound on the obstacle
diameter reduces the infinitely many translates to finitely many close
ones for this verification. The selected chord and its reverse
therefore remain physical ground channels in the periodic billiard.

We verify the all-degree support-to-graph Jacobian, including its sign.
At the right contact, the stationary support-envelope formula is
$x(y)=(h(\theta)-y\sin\theta)/\cos\theta$. At fixed $y$,
$\partial_s x=\partial_s h(\theta(y))/\cos\theta(y)$.
Here $\sin\theta(y)=\kappa_0y+O(y^2)$ and the inward graph is
$1-x(y)$. At the left contact use longitudinal coordinate $-x$ while
retaining global transverse coordinate $y$. Equivalently the outward
normal is $(-\cos\phi,\sin\phi)$, $\theta=\pi-\phi$;
again $\sin\theta(y)=\kappa_1y+O(y^2)$ and the inward graph
is one minus the outward longitudinal coordinate. Consequently
\begin{equation}\label{eq:supportjac}
 \frac{\partial q_{b,r}}{\partial s_{br}}=-r!\kappa_b^r,
 \qquad
 \frac{\partial q_{b,r}}{\partial s_{1-b,r}}=0,
                   \qquad 3\le r\le K.
\end{equation}
The off-contact weight vanishes to order two. A degree-$r$ perturbation
changes no graph jet of lower degree. The compensator vanishes to
order $2K+2$ at both contacts and changes none of the displayed jets.
The Jacobian of $s\mapsto(q_{b,r})$ is thus block lower triangular
with invertible diagonal blocks, without assuming evenness.

The obstacle area and its free-area derivative are
\begin{equation}\label{eq:area}
 \mathcal S(h)=\frac12\int_0^{2\pi}(h^2-(h')^2)\dd\theta,
 \qquad
 \partial_\zeta A=-\int_0^{2\pi}(h+h'')\sin^{2K+2}\theta
                                      \dd\theta<0.
\end{equation}
In the order $(q,A)$ the full Jacobian in $(s,\zeta)$ has an
invertible jet block, a zero jet--$\zeta$ column, and the nonzero
last entry \eqref{eq:area}. The analytic inverse-function theorem
gives the claimed free-area coordinates. The analytic implicit-function
theorem also gives $\zeta=\zeta(s,a,b)$ with $\mathcal S(h)=\pi$
near the unit disk, proving the fixed-area assertion for nearby radii.
All leading selected coefficients depend only on $g,A,\kappa_0,\kappa_1$,
as follows by retaining the quadratic stationary action and constant
flux. These are fixed on this last subfamily. The statement is about
the selected hierarchy, not all channels in the table.
\end{proof}

\subsection{A dimension-sharp coordinate map of scalar means}
The following elementary interpolation argument is included to distinguish
an exact-germ inverse from a finite physical design. Fix a compact small
coordinate neighborhood in Proposition~\ref{prop:realization}.
Use $n_e,n_o$ as in Theorem~\ref{thm:main}; then
$n_e+n_o=K-2$. All constants may depend on this fixed $K$ and family.

\begin{proposition}[Positive-window coordinates]\label{prop:windows}
There exist fixed positive offsets and a compact parameter ball with
nonempty interior such that $2K-4$ scalar means give bi-Lipschitz
coordinates when $A$ is fixed, and $2K-3$ do so when $A$ is free.
The fixed-area design uses $n_e$ values of $P_b$ and $n_o$ values
of $R_b$ in each orientation. The free-area design uses $n_e+1$
values of $P_0$, $n_e$ values of $P_1$, and $n_o$ values of each
$R_b$. Empty blocks are omitted. The offsets and their conditioning
constants need not be uniform as $K$ grows.
\end{proposition}
\begin{proof}
For fixed area use the coefficient coordinates $(\xi,\eta)$ in
Theorem~\ref{thm:main}. On the analytic physical family,
\[
 G_b(d)=1+\sum_{k=1}^{n_e}\xi_{b,k}d^k
                   +d^{n_e+1}U_b(d),\qquad
 J_b(d)=\sum_{k=1}^{n_o}\eta_{b,k}d^k
                   +d^{n_o+1}V_b(d),
\]
with uniformly bounded remainders and first parameter derivatives.
For a block of length $n\ge1$ choose nodes $h,2h,\ldots,nh$.
The nonconstant nodal matrix is
$V_{n,h}=((ih)^k)_{1\le i,k\le n}$ and satisfies
$\|V_{n,h}^{-1}\|\le C_nh^{-n}$ for $0<h\le1$.
Preconditioning the Jacobian block by this inverse turns the remainder
into $O(h)$. For sufficiently small fixed $h$ the combined Jacobian
is invertible. Passing to raw means multiplies rows by known nonzero
constants $\alpha(ih)^2$ and restores the known count intercepts.

When area is free, use the analytic coordinates
\begin{equation}\label{eq:scaledcoordinates}
 \alpha,\quad (\alpha\xi_{b,k})_{b,k},\quad
                      (\alpha\eta_{b,k})_{b,k}.
\end{equation}
They are coordinates because $\alpha>0$ and $A\mapsto\alpha$
is invertible. Divide each raw mean by its known offset squared,
not by the unknown area. The polynomial count blocks have one shared
constant coefficient $\alpha$. The first orientation has $n_e+1$
positive nodes, the second $n_e$ positive nodes. Their combined
$2n_e+1$ square matrix is invertible. Indeed, an element of its kernel
gives a degree-$n_e$ polynomial with $n_e+1$ distinct roots in the
first orientation. That polynomial and hence the shared constant are
zero. The other polynomial has its $n_e$ positive roots and the root
zero, so is zero as well. The inverse has norm $O(h^{-n_e})$,
including the case $n_e=0$. The moment blocks have no constant and
are the matrices $V_{n_o,h}$. After block preconditioning each
remainder is $O(h)$ in the full parameter derivative, even when
$n_e$ and $n_o$ differ. A sufficiently small fixed $h$ gives an
invertible Jacobian. No evaluation at zero is part of this design.

At the center multiply the full Jacobian by its inverse. On a sufficiently
small convex parameter ball the resulting derivative stays within
$1/2$ of the identity. Integration along a segment proves a lower
Lipschitz bound, and a derivative bound proves the upper one.
Take the compact closure of this ball. The nodes have already been
chosen inside a common physical collar where every count probability
is strictly between zero and one. Continuity and compactness now give
fixed positive probability margins; no subsequent change of the nodes
is needed. This proves the proposition and Theorem~\ref{thm:main}(iv).
\end{proof}

The numbers above are minimal for differentiable locally bi-Lipschitz
maps from these physical parameter balls into scalar mean coordinates.
Such a map has an injective derivative, so its target dimension is at
least $2K-4$, or $2K-3$ when area varies. This is a dimension statement
for this class of maps, not a universal information-theoretic bound.

\subsection{An explicitly compressed binary experiment}
Choose a supplied constant $T_*>0$ that bounds $|u+v|$ on the fixed
patches at all selected offsets. In a moment channel, record the physical
variable $X=(u+v)\one_E$, use an independent uniform $U$ on $[0,1]$,
and return only
\begin{equation}\label{eq:binary}
 B=\one_{\{U\le1/2+X/(4T_*)\}}.
\end{equation}
The compression seed $U$ and the uncompressed value $X$ are not returned.
This produces a Bernoulli mean $1/2+R_b(d)/(4T_*)$ in
$[1/4,3/4]$. Failure still consumes a preparation and gives $X=0$.
The endpoint sensor is used before compression; this is not a
count-only experiment. For a probability channel simply return
$\one_{E_{2,b}(d)}$.

\begin{theorem}[Fixed-dimensional charged observation]\label{thm:experiment}
On either compact physical parameter ball of Proposition~\ref{prop:windows},
with the finite scalar channels implemented as above, there is a Borel
estimator with error at most $\varepsilon$ and probability at least
$1-\delta$ using
\[
 N_{\mathrm{tot}}\le C\varepsilon^{-2}\log(C/\delta)
\]
fresh preparations, for small $\varepsilon>0$ and $0<\delta<1/4$.
All failures and auxiliary randomizations are included. For at most
$N$ fresh preparations, with adaptive or randomized choice among these
specified binary channels, the squared minimax risk is between $c/N$
and $C/N$ for all sufficiently large $N$. These assertions do not give
a lower bound for the uncompressed, real-valued endpoint record.
\end{theorem}
\begin{proof}
The finite vector of binary means is an invertible affine row
transformation of the mean map in Proposition~\ref{prop:windows}.
It is therefore bi-Lipschitz. Its finitely many Bernoulli probabilities
have positive margins on the compact parameter ball, by the last
paragraph of that proposition and \eqref{eq:binary}.

Use $L$ independent preparations in each of the $s$ fixed channels.
For their sample-mean vector $Y$ and true mean vector $p(\vartheta)$,
the bounded-summand exponential estimate and a union bound give
\[
 \Prob\{\|Y-p(\vartheta)\|_\infty>t\}
                 \le2s\exp(-2Lt^2).
\]
Take a finite ordered net in the parameter ball whose mean image has
mesh at most $t$. Choose the first point minimizing distance to $Y$.
This is a Borel estimator. On the displayed event's complement the
image error is at most a fixed multiple of $t$; the inverse Lipschitz
bound yields parameter error at most $Ct$. Set $t=c\varepsilon$
and $L=\lceil Ct^{-2}\log(2s/\delta)\rceil$.
For squared risk use $L=\lfloor N/s\rfloor$, a net of mesh
$N^{-1/2}$, and the variance bound
$\E\|Y-p(\vartheta)\|^2\le C/L$. This gives the upper bound.

For the lower bound choose two interior physical parameters
$\vartheta_\pm=\vartheta_*\pm aN^{-1/2}v$, with $\|v\|=1$.
For every allowed channel, smoothness and the positive probability
margins imply
\[
 D\{\operatorname{Ber}(p_+)\Vert\operatorname{Ber}(p_-)\}
 \le\frac{(p_+-p_-)^2}{p_-(1-p_-)}\le Ca^2/N.
\]
Condition on the past and on the parameter-independent randomization
used to select the next channel. Do not reveal the compression seed
$U$, which is not part of the observed record.
The channel-selection rule is the same under the two alternatives.
The entropy chain rule bounds the full-record relative entropy by
$Ca^2$; an early stopping rule can be padded with parameter-independent
bits. Choose $a$ so that total variation is at most $1/2$.
Any estimator induces the nearer-parameter test, whose sum of error
probabilities is at least $1/2$. On an error its squared parameter
loss is at least $a^2/N$, yielding a maximum risk at least $a^2/(4N)$.
The alternatives are the actual support-function tables just constructed.
The argument uses only the specified compressed Bernoulli likelihoods,
not a likelihood for the moving-support full endpoint experiment.
\end{proof}

\section{Physical coordinates and intrinsic finite-window experiments}
\label{sec:intrinsicexperiments}
The original Cartesian realization and experiment have been retained in
full. We record the precise extensions, rather than using the word
``count'' for a sensor that first observes a real-valued mark.

\begin{corollary}[Physical leading and higher coordinates]
\label{cor:physicalleading}
At each fixed $K\ge3$ there are analytic periodic families on which
\[
 g,\ \kappa_0,\ \kappa_1,\ A,
       \quad(q_{b,r})_{b=0,1;\,3\le r\le K}
\]
are local coordinates. The calibrated intrinsic finite-jet map of
Corollary~\ref{cor:leadingchart} is a coordinate map on these physical
families, of dimension $2K$ on a chosen orientation cover.
\end{corollary}
\begin{proof}
Use the support function \eqref{eq:support}, allow its two radii
$R_0,R_1$ to vary near one, and replace the horizontal lattice period
$3$ by $2+g$, with $g$ near one; keep the vertical period four.
The selected support points remain $(1,0)$ and $(-1,0)$, so the gap
is $g$ and the curvatures are $R_b^{-1}$. The jet diagonal
\eqref{eq:supportjac} is unchanged. The area compensator has derivative
\eqref{eq:area}, still nonzero, while
$A=4(2+g)-\mathcal S(h)$. In the variable order consisting of gap,
radii, support jets and compensator, the full coordinate Jacobian has
invertible leading, jet and area diagonal blocks. The analytic
inverse-function theorem applies. Positivity and clearance persist
near the original table. Composition with the calibrated inverse proves
the asserted intrinsic coordinate statement.
\end{proof}

\begin{corollary}[Intrinsic positive-window observation]
\label{cor:arcwindows}
On the fixed-leading-geometry families of
Proposition~\ref{prop:realization}, every assertion of
Proposition~\ref{prop:windows} remains true after replacing the
Cartesian means $R_b$ by the intrinsic means $M_b$. The charged
binary experiment of Theorem~\ref{thm:experiment} also remains true,
with the bounded mark
$X=(s_-+s_+)\one_E$ and a fixed bound $|X|\le T_*$.
Its risk bounds concern that specified compressed experiment only.
\end{corollary}
\begin{proof}
Proposition~\ref{prop:arcinverse} replaces the coefficient coordinates
$(\xi,\eta)$ by $(\xi,\widetilde\eta)$ without changing their
triangular invertibility. On the analytic physical families all finite
remainders and their first parameter derivatives are uniformly bounded.
The Vandermonde and shared-intercept argument in
Proposition~\ref{prop:windows} therefore applies verbatim, including
its raw-probability normalization when area is free. The arclength mark
is bounded on the fixed compact patches. Returning only
$\one_{\{U\le1/2+X/(4T_*)\}}$, with the seed and original mark hidden,
gives mean $1/2+M_b/(4T_*)$. All margin, entropy and estimation arguments
in Theorem~\ref{thm:experiment} then apply with the same stated scope.
\end{proof}

Neither corollary claims finite-sample stability of analytic continuation
or lattice-network reconstruction. In particular the onset $T_0$ and
quadratic limits in the exact-germ calibration are not treated as
noiseless outputs of a finite positive-window experiment. The statistical
rate above is for a fixed-dimensional, fixed-leading-geometry family
and a specified finite collection of binary channels.

\section{Relative laws, normal form, and the retained smooth invariant}
\label{sec:relative}
The asymmetric inverse above is a short-record theorem. It neither
requires nor replaces the function-valued long-record result. We state
the distinction precisely and include the analytic comparison in a
form with an explicit physical chart convention.

\subsection{The smooth physical law}
For $j$ alternating flights starting at contact $0$, write
$E_j=W_j-jg$, $p=j\bmod2$, and
$b_j=-W_{j,uv}/d_j^0$. Define
\[
 a_b=\frac{\sqrt{c_0c_1}\sinh\gamma}{gc_{1-b}},\qquad
 q_p=\sqrt{a_0a_p},\qquad d_j^0=\frac{q_p}{\sinh(j\gamma)}.
\]
The complete smooth result proved in \cite{Companion}, in the sections
on the half-line boundary layers and differentiated operators, is as
follows. On fixed common neighborhoods there are positive smooth $B_b$
and strictly convex smooth $S_b$ with
\[
 B_b(0)=1,\quad S_b(0)=S'_b(0)=0,\quad S_b''(0)=a_b,
\]
such that for each fixed derivative order $k$,
\begin{equation}\label{eq:smoothrelative}
 \|E_j-S_0\oplus S_p\|_{C^k}
       +\|b_j-B_0\otimes B_p\|_{C^k}\le C_k\tau^j,
                    \qquad 0<\tau<1.
\end{equation}
The norms include fixed mixed geometric derivatives on the compact
local families specified there. Normalization of the exact cofactor
identity precedes trace-norm comparison of the two end perturbations;
an uncontrolled absolute error is never divided by $d_j^0$.
The common Morse domain then gives, down to the onset,
\begin{equation}\label{eq:integratedrelative}
 \left\|\frac{2A\sinh(j\gamma)}{d^2}P_{j,0}(d)
       -\mathcal F_{0p}(d)\right\|_{C^k([0,d_*])}
                              \le C_k\tau^j,
\end{equation}
where
\begin{equation}\label{eq:limitlaw}
 \mathcal F_{bc}(d)=\frac{\sqrt{a_ba_c}}{\pi d^2}
       \int(d-S_b(u)-S_c(v))_+B_b(u)B_c(v)\dd u\dd v.
\end{equation}
The corresponding result holds with the other starting contact.
These are smooth, parameter-differentiated statements on a common
physical collar, not merely Taylor-series identities.

The full proofs and all earlier applications remain in the companion
volume, including the first-hit localization, common-domain integration,
operator estimates, and both parities. The present theorem does not
use an asymptotic approximation as a substitute for the exact
finite-flight density: Lemma~\ref{lem:flux} supplied that density
directly.

\subsection{The analytic comparator in physical coordinates}
Suppose analytic canonical coordinates put a local return map in the
form
\begin{equation}\label{eq:normalform}
 N(s,p)=(\Delta(sp)s,\Delta(sp)^{-1}p),\quad
                   \Delta(0)=\lambda\in(0,1).
\end{equation}
This is the local hyperbolic normal-form setting recalled in
\cite[Section 2]{DSKL}; its local use is distinct from the hypotheses
of that paper's global inverse theorem. In this subsection the charts
are supplied, not constructed for an arbitrary smooth family.

\begin{proposition}[Relative projection formula]\label{prop:normalform}
For the mixed endpoints $s_0=s$, $p_n=t$, the small solution satisfies
\begin{equation}\label{eq:mixed}
 I=st\Delta(I)^n,\quad p_0=t\Delta(I)^n,\quad
 s_n=s\Delta(I)^n,\quad
 T_n:=\partial_t p_0=
 \frac{\Delta(I)^n}{1-nI\Delta'(I)/\Delta(I)}.
\end{equation}
On a fixed smaller neighborhood, for every fixed $k$ and some
$\rho\in(\lambda,1)$,
$\|\lambda^{-n}T_n-1\|_{C^k}\le C_k\rho^n$.
Let the physical canonical charts be
\[
 (u,P)=(U(s,p),\mathcal R(s,p)),\quad
 (v,Q)=(V(q,t),\mathcal Z(q,t)),\quad U_s(0)V_t(0)\ne0,
\]
with generating convention $P=-W_{n,u}$, $Q=W_{n,v}$.
For $\Phi_n(s,t)=(U(s,p_0),V(s_n,t))$ and
$\mathcal D_n=\det D\Phi_n$, the exact physical twist is
\begin{equation}\label{eq:projection}
 |W_{n,uv}|=\left|\frac{T_n}{\mathcal D_n}\right|
                       \circ\Phi_n^{-1},\qquad
 \mathcal D_n(0)=U_s(0)V_t(0)-\lambda^{2n}U_p(0)V_q(0).
\end{equation}
The inverses exist on one physical box for all sufficiently large $n$.
If $\sigma(u)$ and $\theta(v)$ invert $U(s,0)$ and $V(0,t)$,
the normalized physical twist converges there in every fixed $C^k$
norm to
\begin{equation}\label{eq:projectionfactors}
 \frac{U_s(0)}{U_s(\sigma(u),0)}
             \frac{V_t(0)}{V_t(0,\theta(v))}.
\end{equation}
\end{proposition}
\begin{proof}
Choose $\lambda<\mu<\rho<1$ and a complex disk where $\Delta$
is holomorphic, nonzero, and bounded by $\mu$. On a sufficiently small
endpoint bidisk the map $I\mapsto st\Delta(I)^n$ is a contraction
uniformly in $n$, since $\sup_n n\mu^{n-1}<\infty$.
It yields a holomorphic solution $I=O(\mu^n)$. Invariance of $sp$
gives the middle two formulas in \eqref{eq:mixed}; the intermediate
points $(s\Delta(I)^i,t\Delta(I)^{n-i})$ stay in the chart.
Implicit differentiation gives its last formula, valid also at $s=0$.
The analytic logarithm gives
\[
 \Delta(I)^n/\lambda^n=1+O(n\mu^n),\qquad
 1-nI\Delta'(I)/\Delta(I)=1+O(n\mu^n).
\]
These are estimates of normalized quantities. Cauchy bounds on a
smaller fixed bidisk and $n\mu^n\le C\rho^n$ give every fixed
derivative estimate.

It follows that $\Phi_n$ tends in each fixed derivative norm to
$(U(s,0),V(0,t))$. Transversality, a smaller physical box, and the
inverse-function contraction give inverse maps on that same box,
converging with their derivatives. Pulling back
$\dd u\wedge\dd P=\dd s\wedge\dd p_0$ to $(s,t)$ gives
$(-W_{n,uv})\mathcal D_n=T_n$. Its absolute value is
\eqref{eq:projection}; the derivative at zero gives the determinant
there. Dividing by the exact origin value and using convergence of the
inverse maps gives \eqref{eq:projectionfactors}. Each scalar derivative
has constant sign on the chosen box, so its normalized ratio is positive.
\end{proof}

The same-section specialization needs an identification of charts, not
just an equality of section names. Choose the same canonical chart at
both ends with $(q,t)$ in the same stable--unstable order as $(s,p)$,
and the same physical transverse and momentum conventions. Then
$V(q,t)=U(q,t)$, $V_t(0)=U_p(0)$, and $V_q(0)=U_s(0)$.
In particular $U_s(0)U_p(0)\ne0$, and the invariant formula becomes
\begin{equation}\label{eq:samesection}
 |W_{n,uv}(0)|=
 \frac{\lambda^n}{|U_s(0)U_p(0)|(1-\lambda^{2n})}.
\end{equation}
A reversal or interchange of a chart must instead be substituted into
\eqref{eq:projection}. For the two-step billiard return, $j=2n$ and
$\lambda=e^{-2\gamma}$. This recovers the exact hyperbolic-sine
normalization rather than only its leading exponential.

Applying the analytic comparison and \eqref{eq:smoothrelative} to the
same analytic physical problem identifies their normalized endpoint
factors by uniqueness of the limit and evaluation on the two axes.
Thus the analytic product phenomenon is not itself the new inverse
mechanism. The odd blocks in \eqref{eq:oddblock} use physical moment
information that the normal form by itself does not supply.

\subsection{The information retained after transverse symmetrization}
Define a locally finite measure on $[0,d_*]$ by
\[
 \nu_b=(S_b)_*\{\sqrt{a_b}B_b(u)\dd u\}.
\]
Near zero its density has the form
$\sqrt2\,x^{-1/2}\mathcal V_b(x)$, where
$\mathcal V_b$ is smooth and $\mathcal V_b(0)=1$.
The two inverse branches of the strictly convex $S_b$ are added in
this density. If $H_{bc}(d)=d^2\mathcal F_{bc}(d)$, then
\[
 H_{bc}(d)=\frac1\pi\int(d-x-y)_+\nu_b(\dd x)\nu_c(\dd y).
\]
Hence the smooth count law is a symmetrized energy pushforward.
The Volterra uniqueness and compatibility proof in \cite{Companion}
recovers the entire $\mathcal V_0,\mathcal V_1$ from the two
diagonal laws and predicts the mixed law. It is not an identification
of two smooth functions from their Taylor series.

The same volume gives the weighted Abel inverse norm, the two-derivative
gain from integrated flux, and the separately charged smooth-class
calibration procedure. Its full-profile preparation bound is sufficient,
not a matching minimax theorem. The limiting even-contact inverse and
its finite-jet equivalence to probability-only two-flight coefficients
remain valid with their original hypotheses. They are distinct from
the all-degree signed inverse proved here. In particular a classical
analytic normal form, a smooth symmetrized profile, and an asymmetric
contact graph are three different mathematical objects.

\subsection*{Supplementary material}
Supplement S, \emph{Nonlinear boundary laws: complete geometric proofs
and auxiliary experiments}, is supplied as part of this manuscript's
submission package, not cited as an independently published or accepted
paper. It contains the full smooth relative-law, Volterra, Abel and
auxiliary-experiment proofs summarized in Section~\ref{sec:relative}.
The preceding primary proof of intrinsic and registered rigidity is
self-contained and uses none of those supplementary results as a lemma.
The preceding Cartesian proofs are retained in full; the earlier primary
article is additionally preserved as an archival source, not as a second
new submission. The source-distribution index identifies the respective
volumes and their immutable provenance.
\subsection*{Acknowledgment and assisted analysis}
AI-assisted analysis contributed to the intrinsic arclength calibration,
the filtered inverse, the registered-network reconstruction, and the
count-jet fiber argument. Earlier signed-moment calculations and the
complete smooth theory are retained with their original provenance.
The named author is responsible for checking the arguments and references
before submission. Finite arithmetic diagnostics and compilation are
reproducibility evidence, not formal proof certificates.
\begin{thebibliography}{99}
\raggedright
\bibitem{BDKL}
P. B\'alint, J. De Simoi, V. Kaloshin, and M. Leguil,
\emph{Marked length spectrum, homoclinic orbits and the geometry of open
dispersing billiards}, Comm. Math. Phys. \textbf{374} (2020), 1535--1575.
\href{https://doi.org/10.1007/s00220-019-03448-x}{doi:10.1007/s00220-019-03448-x}.
\bibitem{DSKL}
J. De Simoi, V. Kaloshin, and M. Leguil,
\emph{Marked length spectral determination of analytic chaotic billiards
with axial symmetries}, Invent. Math. \textbf{233} (2023), 829--901;
arXiv:1905.00890v4.
\href{https://doi.org/10.1007/s00222-023-01191-8}{doi:10.1007/s00222-023-01191-8}.
\bibitem{DSKW}
J. De Simoi, V. Kaloshin, and Q. Wei,
\emph{Dynamical spectral rigidity among $\mathbb Z_2$-symmetric strictly
convex domains close to a circle}, with Appendix B coauthored by H. Hezari,
Ann. of Math. (2) \textbf{186} (2017), 277--314.
\href{https://doi.org/10.4007/annals.2017.186.1.7}{doi:10.4007/annals.2017.186.1.7}.
\bibitem{FL}
D. Finamore and M. Leguil,
\emph{A CAT(0)-approach to the marked length spectral rigidity of Sinai
billiards}, arXiv:2510.18983v1 (2025).
\bibitem{Koval}
I. Koval, \emph{Billiard tables with analytic Birkhoff normal form are
generically Gevrey divergent}, arXiv:2403.14448v1 (2024).
\bibitem{Osterman}
O. V. Osterman, \emph{On length spectrum rigidity of dispersing billiard
systems}, arXiv:2208.12244v1 (2022).
\bibitem{Companion}
Q. Qi, \emph{Nonlinear boundary laws: complete geometric proofs and
auxiliary experiments}, Supplement S to the present submission (2026).
Included supplementary manuscript; not an independently published article.
\end{thebibliography}

\end{document}
